\cotefile{76}
\scanpage{19}
% L1–4
On trouve que $n\longmapsto D[n]$ devient un ens quasi-simplicial, i.e un \unsure{contrafoncteur} sur la catégorie $\unsure{\mathbb{\Delta}|\Delta}$ des ens finis non vides de la forme $\Delta_n=[0,n]_{\mathbb N}$, avec comme
% L5–7
morphismes les appl. \uline{injectives}. Cet ens quasi-simplicial satisfait la condition que
% L8
$\forall n,m\in\mathbb N$ avec $n<m$, et $d_m\in D[m]$, l'application
% L9–10
\[
\begin{aligned}
\mathrm{Hom}_{\mathbb{\Delta}}(\Delta_n,\Delta_m)&\longrightarrow D_n\\
u&\longmapsto u^*(d_m)
\end{aligned}\tag*{\unsure{(*)}}
\]
% L11–12
est injective, qui signifie aussi que la catégorie
% L13
\[
\widehat{\mathbb{\Delta}}/D[\,]
\]
% L14
est ordonnée (\unsure{ou mieux|ou même|ou au moins|i.e. équivalente}, préordonnée).

% L15–17
L'ens ordonné $D$ se reconstitue en \deleted{\unsure{fon|fonction}} \unsure{termes|fonction} du $D[\,]$, comme l'ens ordonné \deleted{\unsure{que}} défini par la dite catégorie ordonnée. \deleted{\unsure{Inversement}}, \inserted{\unsure{On trouve}}
% L18
\deleted{la donnée de l'ens \unsure{f}}
% L19–22
\diagram{\begin{tikzcd}[ampersand replacement=\&,column sep=small]
\begin{gathered}\text{géométrie des drapeaux }D\\
\text{avec syst. d'ordres }\text{\inserted{totaux}}\ \unsure{\text{coh.}|\text{etc.}}\ \text{sur}\\
\text{les simplexes des } \unsure{\mathscr D|\Phi}\ (\text{ens des}\\
\text{él. minimaux de }D)\end{gathered}
\arrow[r,leftrightarrow] \&
\begin{gathered}\text{ens quasi-}\\\text{simpliciaux}\\
\text{satisfaisant la}\\\text{condition }(*)\text{.}\end{gathered}
\end{tikzcd}}
% archive sheet number: 18
